% Part: lambda-calculus
% Chapter: introduction
% Section: free-variables

\documentclass[../../../include/open-logic-section]{subfiles}

\begin{document}

\olfileid{lam}{syn}{fv}
\olsection{స్వేచ్ఛా చరాలు}

లాంబ్డా కలనశాస్త్రం ప్రమేయాలకు సంబంధించినది; లాంబ్డా అమూర్తీకరణ
ద్వారానే ఇక్కడ ప్రమేయాలు ఏర్పడతాయి. సహజార్థంలో, $\lambd[x][M]$
అనేది ఆ ప్రమేయపు ఆర్గ్యుమెంటును~$x$కు కేటాయించినప్పుడు~$M$ ఇచ్చే
విలువలు గల ప్రమేయం. కానీ $M$లోని ప్రతి $x$ ఘటనకూ దీనితో సంబంధం
ఉండదు: $M$లో మరో అమూర్తీకరణ $\lambd[x][N]$ ఉంటే, $N$లోని $x$
ఘటనలు $\lambd[x][N]$కు సంబంధించినవి; $\lambd[x][M]$కు కావు.
కాబట్టి $\lambd[x][M]$ లోపలున్న లాంబ్డా అమూర్తీకరణ $\lambd[x]$,
$M$లోని $x$ ఘటనల్లో మరొక లాంబ్డా అమూర్తీకరణ ఇప్పటికే బంధించని
వాటిని---అంటే $M$లోని $x$ \emph{స్వేచ్ఛా} ఘటనలను---\emph{బంధిస్తుంది}.

\begin{defn}[పరిధి]
పదం~$N$లో $\lambd[x][M]$ సంభవిస్తే, దాని శరీరమైన $M$ యొక్క అనుగుణ
ఘటనే ఆ $\lambd[x]$ యొక్క \emph{పరిధి}.
\sourcecorrection{OLTELAMSYN-003}{ఆంగ్ల మూలం $N$లో సంభవించే $\lambd[x][M]$ యొక్క పరిధిగా $N$ యొక్క అనుగుణ ఘటనను పేర్కొంది. అమూర్తీకరణ యొక్క పరిధి దాని శరీరమైన $M$; తదుపరి ఉదాహరణలూ దానినే నిర్ధారిస్తున్నాయి. ఆ స్థానంలో $M$ను పునరుద్ధరించాం.}
\end{defn}


\begin{defn}[స్వేచ్ఛా, బద్ధ ఘటన]
  పదం~$M$లో చరం~$x$ యొక్క ఒక ఘటన ఏ $\lambd[x]$ పరిధిలోనూ
  లేకపోతే అది \emph{స్వేచ్ఛా} ఘటన; లేకపోతే \emph{బద్ధ} ఘటన.
  $\lambd[x][M]$లో చరం~$x$ యొక్క ఒక ఘటనను మొదటి $\lambd[x]$
  బంధించేది, $M$లోని ఆ $x$ ఘటన స్వేచ్ఛగా ఉన్నప్పుడు, అప్పుడు మాత్రమే.
\end{defn}

\begin{ex}
  $\lambd[x][x y]$లో $x$, $y$ రెండూ $\lambd[x]$ పరిధిలో ఉన్నాయి;
  అందువల్ల $x$ను $\lambd[x]$ బంధిస్తుంది. $y$ ఏ $\lambd[y]$
  పరిధిలోనూ లేదు కనుక అది స్వేచ్ఛగా ఉంది. $\lambd[x][x x]$లోని
  రెండు $x$ ఘటనలూ $\lambd[x]$చే బద్ధమైనవి; ఎందుకంటే $xx$లో
  రెండూ స్వేచ్ఛగా ఉన్నాయి. $((\lambd[x][xx])x)$లో చివరి $x$
  స్వేచ్ఛగా ఉంది; అది ఏ $\lambd[x]$ పరిధిలోనూ లేదు.
  $\lambd[x][(\lambd[x][x])x]$లో మొదటి $\lambd[x]$ పరిధి
  $(\lambd[x][x])x$; రెండవ $\lambd[x]$ పరిధి చివరి నుంచి రెండవ
  $x$ ఘటన. $(\lambd[x][x])x$లో చివరి $x$ ఘటన స్వేచ్ఛగా ఉంది;
  చివరి నుంచి రెండవది బద్ధంగా ఉంది. కనుక
  $\lambd[x][(\lambd[x][x])x]$లో చివరి నుంచి రెండవ $x$ ఘటనను
  రెండవ $\lambd[x]$ బంధిస్తుంది; చివరి ఘటనను మొదటి $\lambd[x]$
  బంధిస్తుంది.
\end{ex}

పదం~$P$లోని అన్ని చర ఘటనలను పరిశీలిస్తే, స్వేచ్ఛా చరాల సమితి
లభిస్తుంది. ఆ సమితిని $\FV{P}$తో సూచిస్తాం; దానికి సహజమైన
నిర్వచనం ఇది:

\begin{defn}[పదం యొక్క స్వేచ్ఛా చరాలు] \ollabel{def:fv}
  పదం యొక్క \emph{స్వేచ్ఛా చరాల} సమితిని ఆగమన పద్ధతిలో ఇలా
  నిర్వచిస్తాం:
  \begin{enumerate}
    \item $\FV{x} = \{x\}$ \ollabel{def:fv1}
    \item $\FV{\lambd[x][N]} = \FV{N} \setminus \{x\}$ \ollabel{def:fv2}
    \item $\FV{PQ} = \FV{P} \cup \FV{Q}$ \ollabel{def:fv3}
  \end{enumerate}
\end{defn}

\begin{prob}
  \begin{enumerate}
  \item $\lambd[g][(\lambd[x][g (x x)]) \lambd[x][g (x x)]]$
    అనే పదంలోని $\lambd[g]$, రెండు $\lambd[x]$ల పరిధులను
    గుర్తించండి.
  \item $\lambd[g][(\lambd[x][g (x x)]) \lambd[x][g (x x)]]$లో
    చరాల ఘటనలన్నీ బద్ధమైనవేనా? ఆయా ఘటనలను ఏ అమూర్తీకరణలు
    బంధిస్తాయి?
  \item $\FV{\lambd[x][(\lambd[y][(\lambd[z][x y]) z]) y]}$
    విలువను ఇవ్వండి.
  \end{enumerate}
\end{prob}

\begin{explain}
స్వేచ్ఛా చరం బయటి ప్రపంచానికి (\emph{పరిసరానికి}) ఉన్న సూచన
వంటిది. స్వేచ్ఛా చరాలు గల పదాన్ని కొంతవరకే నిర్దేశించిన పదంగా
చూడవచ్చు; ఎందుకంటే దాని ప్రవర్తన మనం పరిసరాన్ని ఎలా ఏర్పాటు
చేస్తామనే దానిపై ఆధారపడుతుంది. ఉదాహరణకు, ఆర్గ్యుమెంట్~$x$ను
తీసుకుని దానిపై $f$ను ప్రయోగించే $\lambd[x][f x]$లో చరం~$f$
స్వేచ్ఛగా ఉంది. ఆ పదం విలువ, అది ఉన్న పరిసరంపై, ముఖ్యంగా ఆ
పరిసరంలో $f$ విలువపై ఆధారపడుతుంది.

ఈ పదంపై అమూర్తీకరణ ప్రయోగిస్తే $\lambd[f][\lambd[x][f
    x]]$ వస్తుంది. ఈ పదం ఇక పరిసరంలోని $f$ చరంపై ఆధారపడదు;
ఎందుకంటే ఇది రెండు ఆర్గ్యుమెంట్లను తీసుకుని మొదటిదాన్ని రెండవదానిపై
ప్రయోగించిన ఫలితాన్ని ఇచ్చే ప్రమేయాన్ని సూచిస్తుంది. పరిసరంలో $f$
విలువను మార్చినా ఈ పదం ప్రవర్తనపై ప్రభావం ఉండదు: ఇది ఆర్గ్యుమెంటుగా
అందినదానినే వాడుతుంది; పరిసరంలోని $f$ విలువను కాదు.
\end{explain}

\begin{defn}[సంవృత పదం, సంయోజకం]
  స్వేచ్ఛా చరాలు లేని పదాన్ని \emph{సంవృత పదం} లేదా
  \emph{సంయోజకం} అంటాం.
\end{defn}

\begin{lem}\ollabel{lem:fv}
  \begin{enumerate}
    \item \ollabel{lem:fv-abs} $y \neq x$ అయితే, $y \in
      \FV{\lambd[x][N]}$ అయ్యేది, $y \in \FV{N}$ అయినప్పుడు,
      అప్పుడు మాత్రమే.
    \item \ollabel{lem:fv-app} $y \in \FV{PQ}$ అయ్యేది,
      $y \in \FV{P}$ లేదా $y \in \FV{Q}$ అయినప్పుడు,
      అప్పుడు మాత్రమే.
    \end{enumerate}
\end{lem}

\begin{proof}
  సాధన.
\end{proof}

\begin{prob}
  \olref[lam][syn][fv]{lem:fv}ను నిరూపించండి.
\end{prob}

\end{document}
